\documentclass[11pt]{article}

\usepackage[margin=1in]{geometry}
\usepackage{amsmath,amssymb,amsthm}
\usepackage{booktabs}
\usepackage{hyperref}
\usepackage{natbib}

\title{Do Biological Abstractions Generalize to Meta-Level AI Systems Engineering?\\
\large Evidence from Evolutionary Configuration Search in Operon}

\author{Bogdan Banu}

\date{April 2026}

\begin{document}

\maketitle

\begin{abstract}
We test whether biological abstractions designed for running multi-agent
organisms also generalize to \emph{evolving} their configurations.
Using Operon's biomimetic framework, we implement an evolution loop that
maps candidate configurations to Genome objects, detects stagnation via
an extended EpiplexityMonitor, and compares tournament mutation against
LLM-backed proposal with rich filesystem context.

Our findings are mixed: biological abstractions generalize cleanly as
\textbf{code structure} (lossless Genome round-trip in 5~lines,
scale-invariant epistemic health monitoring) but \textbf{not} as
optimization algorithms (tournament mutation matches or outperforms LLM
reasoning across config-space and topology-space search). Rich context
improves the LLM proposer $3\times$ over compressed history, consistent
with the Meta-Harness insight, but does not achieve dominance over blind
mutation. We relate our architecture to the categorical framework of
de~los~Riscos et~al.\ and argue that the structural guarantee features
(immune systems, developmental gating, epistemic monitoring) remain the
core value proposition of biology-inspired agent engineering.
\end{abstract}

\section{Introduction}

Biology-inspired abstractions have been proposed for structuring AI agent
systems~\cite{banu2026operon}. Operon implements genes, genomes,
epiplexity monitoring, immune systems, and developmental gating as
engineering primitives for multi-agent organisms. Phases C1--C7 validated
these abstractions at the organism level: adapters bridge six external
frameworks, four TLA+ specifications verify safety invariants, and
compilers translate organisms across deployment targets.

Phase~C8 asks a different question: do these abstractions generalize to
the \emph{meta-level}---evolving organisms, not just running them? This
is the natural extension of the biological metaphor: if organisms have
genomes, can we evolve those genomes to discover better configurations?

We test this by implementing a meta-harness that:
\begin{enumerate}
\item Maps \texttt{CandidateConfig} objects to \texttt{Genome} via the
    existing gene abstraction
\item Uses \texttt{EpiplexityMonitor} with a pluggable
    \texttt{DistanceProvider} for stall detection
\item Wraps each evolution step as a Zardini \texttt{DesignProblem}
\item Compares tournament mutation against an LLM proposer with
    filesystem-backed context (following the Meta-Harness
    insight~\cite{lee2026metaharness})
\end{enumerate}

\section{Related Work}

Ao et~al.~\cite{ao2026delegation} prove that without exogenous signals,
delegated multi-agent networks are dominated by centralized baselines.
Lee et~al.~\cite{lee2026metaharness} show that giving proposers
filesystem access to all prior candidates (${\sim}10$~MTok/iteration)
dramatically outperforms compressed-feedback optimizers. de~los~Riscos
et~al.~\cite{delosriscos2026categorical} provide a category-theoretic
framework (\textsc{ArchAgents}) where architectures are objects,
translations are morphisms, and agents are monoidal functors---mapping
tightly to Operon's organism/compiler/adapter ecosystem.

\section{Method}

\subsection{Genome Mapping}

Each \texttt{CandidateConfig} (stage modes, models, thresholds) is
flattened into \texttt{Gene} objects with \texttt{GeneType.STRUCTURAL}
for mode/model and \texttt{GeneType.REGULATORY} for visibility flags.
The round-trip \texttt{candidate\_to\_genome} $\rightarrow$
\texttt{genome\_to\_candidate} is lossless---5~lines of flattening
logic, no type coercion. This is the strongest evidence that the gene
abstraction generalizes beyond its original per-agent configuration
use case.

\subsection{Epistemic Health Monitoring}

The existing \texttt{EpiplexityMonitor} (Bayesian surprise via embedding
cosine similarity) is extended with a \texttt{DistanceProvider} protocol.
For configuration space, \texttt{ConfigHammingDistance} measures
field-level mismatches. The core formula ($\hat{E}_t = \alpha \cdot
\text{novelty} + (1-\alpha) \cdot \text{perplexity}$) is unchanged---only
the novelty source is pluggable. The monitor triggers
\textsc{Stagnant}/\textsc{Exploring} transitions identically to the
embedding path, confirming scale-invariance.

\subsection{Dual Stall Detection}

Two independent signals trigger the LLM proposer:
\begin{enumerate}
\item \textbf{Config novelty stall} (EpiplexityMonitor): consecutive
    similar configurations
\item \textbf{Score plateau}: best score not improved in
    $\text{threshold} \times |\text{tasks}|$ steps
\end{enumerate}
The second signal was necessary because tournament mutations always
produce novel configs (high novelty) even when scores stagnate---the
epiplexity-only check never fired.

\subsection{Proposer Strategies}

\textbf{Tournament mutation} (70\%): select top-$k$, mutate one random
field via \texttt{Genome.mutate()}. The mutation handles discrete values
naturally---the gap in \texttt{Genome.replicate()} (numeric jitter only)
does not affect explicit targeted mutations.

\textbf{LLM proposer} (30\%, on stall): reads full candidate configs +
execution trace metadata from filesystem store. Gemini-2.5-flash with
4096~token budget (truncation at 2000 was the hidden bug that masked all
early results).

\subsection{Topology Mutations (Phase~B)}

Extended \texttt{CandidateConfig} with \texttt{edges} (directed wiring
pairs) and three mutation operators: \texttt{add\_stage},
\texttt{remove\_stage}, \texttt{rewire}. DAG execution via
\texttt{WiringDiagram} + \texttt{ResourceAwareExecutor} (parallel groups
via \texttt{ThreadPoolExecutor}).

\section{Results}

\subsection{Phase~A: Configuration Search}

\begin{table}[h]
\centering
\begin{tabular}{lcc}
\toprule
\textbf{Proposer} & \textbf{Mean Score} & \textbf{N} \\
\midrule
Tournament (compressed context) & 0.32 & 38 \\
LLM (compressed context) & 0.15 & 2 \\
Tournament (rich context) & 0.44 & 9 \\
LLM (rich context) & 0.49 & 24 \\
\bottomrule
\end{tabular}
\caption{Phase~A: config-only evolution. Rich context improves LLM
$3\times$ but only marginally exceeds tournament.}
\end{table}

\subsection{Phase~B: Topology Search}

\begin{table}[h]
\centering
\begin{tabular}{lcc}
\toprule
\textbf{Proposer} & \textbf{Mean Score} & \textbf{N} \\
\midrule
Tournament (topology mutations) & 0.60 & 9 \\
LLM (topology mutations) & 0.36 & 30 \\
\bottomrule
\end{tabular}
\caption{Phase~B: topology + config evolution. Tournament \emph{improves}
with topology; LLM \emph{degrades}. Blind mutation handles structural
changes more productively than reasoned proposals.}
\end{table}

\subsection{Abstraction Quality}

\begin{table}[h]
\centering
\begin{tabular}{lcc}
\toprule
\textbf{Abstraction} & \textbf{Generalizes?} & \textbf{Evidence} \\
\midrule
Genome $\rightarrow$ CandidateConfig & Yes & Lossless round-trip, 5 lines \\
EpiplexityMonitor + DistanceProvider & Yes & Scale-invariant transitions \\
DesignProblem wrapping & Yes & Natural composition \\
\texttt{feedback\_fixed\_point} & No & Evolution $\neq$ convergent iteration \\
TrustRegistry & No & Overkill for 2 proposers \\
\bottomrule
\end{tabular}
\caption{Biological abstraction generalization results.}
\end{table}

\section{Discussion}

\subsection{Code Structure vs.\ Optimization}

The central finding: biological abstractions generalize as \emph{code
structure} but not as \emph{optimization algorithms}. The Genome mapping,
EpiplexityMonitor extension, and DesignProblem wrapping produce clean,
composable code. But the evolutionary search itself---mutation, selection,
population management---does not outperform random tournament mutation.

This is consistent with Ao et~al.: without genuinely new exogenous
signals, the multi-agent (LLM proposer) approach cannot beat the
single-agent (tournament) baseline. Rich context helps ($3\times$), but
the search space (config knobs) is too small for LLM reasoning to
provide structural advantage over random exploration.

\subsection{Categorical Interpretation}

In the de~los~Riscos framework, Phase~A explores agents within a fixed
architecture (knowledge-layer changes: $\text{Know}_A$), while Phase~B
explores architecture morphisms (moving between objects in
\textsc{ArchAgents}). Our finding that topology mutations improve
tournament but degrade LLM suggests that the categorical structure
constrains useful mutations---random morphisms that happen to preserve
feasibility are more productive than reasoned morphisms that optimize
inert structure.

\subsection{Implications}

The structural guarantee features---immune systems, epiplexity monitoring,
developmental gating---remain the core value proposition. C8 confirms
they should be the focus of empirical validation, not meta-optimization.
The evolution loop infrastructure is useful as an evaluation tool (moved
to \texttt{eval/meta/}) but does not belong in the structural guarantee
library.

\section{Conclusion}

Biological abstractions in AI systems engineering generalize to the
meta-level as organizing principles, not as optimization advantages. The
gene abstraction covers configuration space. Epistemic health monitoring
is scale-invariant. Co-design composition works at the meta-level. But
evolution does not beat random mutation for the search spaces we tested.
The right direction for biology-inspired AI engineering is measuring
structural safety guarantees, not optimizing search.

\bibliographystyle{plainnat}
\bibliography{../references}

\end{document}
